\documentclass[11pt,a4paper]{article}
\usepackage[margin=1in]{geometry}
\usepackage{amsmath,amssymb,amsthm}
\usepackage{algorithm}
\usepackage{algpseudocode}
\usepackage{booktabs}
\usepackage{hyperref}

\theoremstyle{definition}
\newtheorem{definition}{Definition}[section]
\newtheorem{theorem}{Theorem}[section]
\newtheorem{lemma}[theorem]{Lemma}
\newtheorem{remark}{Remark}[section]

\title{\textbf{Theoretical Foundations, Energy Invariance, and Computational Efficiency of Root Mean Square Normalization (RMSNorm)}}
\author{Team Scaibu \\ \small Email: \texttt{jaydeep.9664920749@gmail.com} \\ \small Architecture and Core Engine Team}
\date{October 2026}

\begin{document}
\maketitle

\begin{abstract}
This document provides the rigorous mathematical specification, invariance properties, and complexity analysis of Root Mean Square Normalization (RMSNorm). Layer normalization hypothesizes that both mean-centering and variance-scaling contribute to training stability. RMSNorm demonstrates that scaling invariance alone provides equivalent regularization while eliminating mean calculation and additive bias parameters, reducing computational latency and memory access overhead. We formalize the transformation, prove scale invariance, analyze computational efficiency, trace a complete worked numerical example, and discuss numerical edge cases.
\end{abstract}

\section{Notation and Mathematical Preliminaries}

Let $X \in \mathbb{R}^{B \times D}$ denote the activation matrix consisting of $B$ sequence tokens or batch instances with hidden dimension $D$.

\begin{itemize}
    \item $B \in \mathbb{N}_{\ge 1}$: Batch size or sequence length.
    \item $D \in \mathbb{N}_{\ge 1}$: Hidden feature dimension.
    \item $x_i \in \mathbb{R}^D$: Feature vector corresponding to row $i \in \{1, \dots, B\}$.
    \item $\text{RMS}(x_i) \in \mathbb{R}_{> 0}$: Root mean square statistic of activation vector $x_i$.
    \item $\gamma \in \mathbb{R}^D$: Channel-wise learned scaling parameter vector.
    \item $\epsilon > 0$: Denominator stabilization constant ($\epsilon = 10^{-6}$).
    \item $Y \in \mathbb{R}^{B \times D}$: Output normalized activation matrix.
\end{itemize}

\begin{table}[h]
\centering
\caption{Summary of Mathematical Symbols for RMSNorm}
\begin{tabular}{lll}
\toprule
\textbf{Symbol} & \textbf{Type / Domain} & \textbf{Description} \\
\midrule
$X$ & $\mathbb{R}^{B \times D}$ & Activation input matrix \\
$x_i$ & $\mathbb{R}^D$ & Feature vector of token $i$ \\
$\text{RMS}(x_i)$ & $\mathbb{R}_{> 0}$ & Root mean square scalar of row $i$ \\
$Y$ & $\mathbb{R}^{B \times D}$ & Normalized output tensor \\
$\gamma$ & $\mathbb{R}^D$ & Learned multiplicative gain vector \\
$\epsilon$ & $\mathbb{R}_{> 0}$ & Regularization constant ($10^{-6}$) \\
\bottomrule
\end{tabular}
\end{table}

\section{Problem Definition and Core Concepts}

\subsection{Intuitive Meaning}
Layer Normalization scales inputs by standard deviation and shifts by mean. Empirical analysis of transformer activations reveals that the primary stabilizing factor is the scaling invariance of feature magnitude, whereas mean-centering is computationally redundant. RMSNorm scales activations strictly by their root mean square energy, saving memory bandwidth and execution cycles.

\subsection{Formal Mathematical Definition}
For each token vector $x_i \in \mathbb{R}^D$:
\begin{align}
\text{RMS}(x_i) &= \sqrt{\frac{1}{D} \sum_{j=1}^D X_{ij}^2 + \epsilon} \\
\bar{X}_{ij} &= \frac{X_{ij}}{\text{RMS}(x_i)} \\
Y_{ij} &= \gamma_j \bar{X}_{ij}
\end{align}

\section{Algorithm Specification}

\begin{algorithm}[H]
\caption{RMSNorm Forward Transformation}
\label{alg:rms_norm}
\begin{algorithmic}[1]
\Require Input activation matrix $X \in \mathbb{R}^{B \times D}$, gain vector $\gamma \in \mathbb{R}^D$, regularizer $\epsilon > 0$.
\Ensure Normalized output matrix $Y \in \mathbb{R}^{B \times D}$, root mean square vector $r \in \mathbb{R}^B$.
\State $B \gets \text{rows}(X)$, $D \gets \text{cols}(X)$
\State Allocate $Y \in \mathbb{R}^{B \times D}$, $r \in \mathbb{R}^B$
\For{$i \gets 1 \textbf{ to } B$}
    \State $S_i \gets \sum_{j=1}^D X_{ij}^2$
    \State $r_i \gets \sqrt{\frac{S_i}{D} + \epsilon}$
    \For{$j \gets 1 \textbf{ to } D$}
        \State $Y_{ij} \gets \gamma_j \cdot \frac{X_{ij}}{r_i}$
    \EndFor
\EndFor
\State \Return $\{Y, r\}$
\end{algorithmic}
\end{algorithm}

\section{Theoretical Guarantees and Scaling Invariance}

\begin{theorem}[Strict Scaling Invariance of RMSNorm]
For any activation vector $x \in \mathbb{R}^D$ and any positive scalar factor $\alpha > 0$:
\begin{equation}
\text{RMSNorm}(\alpha x) = \text{RMSNorm}(x) \quad (\text{as } \epsilon \to 0)
\end{equation}
\end{theorem}

\begin{proof}
Let $x' = \alpha x$. The scaled root mean square is:
\begin{equation}
\text{RMS}(x') = \sqrt{\frac{1}{D}\sum_{j=1}^D (\alpha x_j)^2} = \sqrt{\frac{\alpha^2}{D}\sum_{j=1}^D x_j^2} = \alpha \sqrt{\frac{1}{D}\sum_{j=1}^D x_j^2} = \alpha \text{RMS}(x)
\end{equation}
Substituting into the normalization formula:
\begin{equation}
Y_j' = \gamma_j \frac{\alpha x_j}{\alpha \text{RMS}(x)} = \gamma_j \frac{x_j}{\text{RMS}(x)} = Y_j
\end{equation}
Thus, RMSNorm guarantees exact scale invariance with respect to input activations and weight scaling.
\end{proof}

\subsection{Limitations}
\begin{itemize}
    \item \textbf{Shift Non-Invariance}: Unlike LayerNorm, RMSNorm is not shift-invariant ($\text{RMSNorm}(x + c\mathbf{1}) \ne \text{RMSNorm}(x)$). If downstream layers require activation centering, bias must be introduced explicitly elsewhere.
\end{itemize}

\section{Complexity Analysis}

\subsection{Time Complexity}
\begin{itemize}
    \item \textbf{Sum of Squares}: Computing $\sum x_j^2$ requires $D$ multiplications and $D-1$ additions per row.
    \item \textbf{Scaling}: Applying $\gamma_j \frac{x_j}{\text{RMS}}$ takes $2D$ operations per row.
    \item \textbf{Total Time Complexity}: $\Theta(B \cdot D)$ FLOPs (requiring $\approx 33\%$ fewer operations than LayerNorm).
\end{itemize}

\subsection{Space Complexity}
\begin{itemize}
    \item \textbf{Auxiliary Memory}: Requires storing $r \in \mathbb{R}^B$, requiring $\mathcal{O}(B)$ auxiliary memory.
\end{itemize}

\section{Exactness and Error Analysis}

The operation is mathematically exact up to floating-point representation. By avoiding mean subtraction, RMSNorm avoids an intermediate memory read/write pass, decreasing rounding errors in memory-bandwidth-bound fused kernel implementations.

\section{Worked Numerical Example}

Consider $B = 1$ token with $D = 4$:
\begin{equation}
X = [[2.0, 4.0, 4.0, 6.0]], \quad \gamma = [1.0, 1.0, 1.0, 1.0], \quad \epsilon = 10^{-6}
\end{equation}

\begin{enumerate}
    \item \textbf{Mean Squared Value}:
    \begin{equation}
    \text{MS} = \frac{2^2 + 4^2 + 4^2 + 6^2}{4} = \frac{4 + 16 + 16 + 36}{4} = \frac{72}{4} = 18.0
    \end{equation}
    \item \textbf{Root Mean Square}:
    \begin{equation}
    \text{RMS} = \sqrt{18.0 + 10^{-6}} = \sqrt{18.000001} \approx 4.2426408
    \end{equation}
    \item \textbf{Normalized Outputs}:
    \begin{align}
    Y_{11} &= 1.0 \times \frac{2.0}{4.2426408} \approx 0.4714045 \\
    Y_{12} &= 1.0 \times \frac{4.0}{4.2426408} \approx 0.9428090 \\
    Y_{13} &= 1.0 \times \frac{4.0}{4.2426408} \approx 0.9428090 \\
    Y_{14} &= 1.0 \times \frac{6.0}{4.2426408} \approx 1.4142136
    \end{align}
\end{enumerate}

\section{Numerical Considerations and Edge Cases}

\begin{itemize}
    \item \textbf{All-Zero Activation Vector}: If $x = \mathbf{0}$, $\sum x_j^2 = 0$. Regularizer $\sqrt{\epsilon} = 10^{-3}$ prevents division by zero and outputs $\mathbf{0}$.
    \item \textbf{Kernel Fusion}: RMSNorm maps directly into single-pass CUDA/Triton warp-reduction kernels, eliminating the two-pass barrier required by LayerNorm.
\end{itemize}

\begin{thebibliography}{9}
\bibitem{zhang2019root} B.~Zhang and R.~Sennrich, ``Root Mean Square Layer Normalization,'' in \emph{Advances in Neural Information Processing Systems (NeurIPS)}, pp.~12360--12371, 2019.
\bibitem{touvron2023llama} H.~Touvron et al., ``LLaMA: Open and Efficient Foundation Language Models,'' \emph{arXiv preprint arXiv:2302.13971}, 2023.
\end{thebibliography}

\end{document}
